\scanpage{16}

% loc: upper question
\uline{Question}\qquad $C$ petite catégorie

$\varphi:C\longrightarrow\mathrm{Ens}$ foncteur exact à droite [$C$ stable par $\varinjlim$ finies] alors $\bar\varphi:(\widehat{C^\circ})^\circ\unsure{\to|{-}}\mathrm{Ens}$ qui prolonge $\varphi$ et commute aux petites $\varprojlim$ (donné aussi par $\bar\varphi(\psi)=\operatorname{Hom}(\psi,\varphi)\quad\psi\in\operatorname{Ob}\underline{\operatorname{Hom}}(C,\mathrm{Ens})$) est-il exact à droite, i.e.

% loc: upper left diagram
\diagram{\begin{tikzcd}[ampersand replacement=\&] {\bar C=(\widehat{C^\circ})^\circ=\underline{\operatorname{Hom}}(C,\mathrm{Ens})^\circ}\arrow[dr]\&\\ C\arrow[u,dash]\arrow[r,"\varphi"]\&\mathrm{Ens}\end{tikzcd}}

% loc: first functor
\[\psi\longmapsto\operatorname{Hom}(\psi,\varphi)\]
transforme-il lim. projectives finies en lim inductives ?

% loc: upper NB
\uline{NB}\quad Par hypoth. sur $\varphi$, c'est le cas pour les $\underset{\longleftrightarrow}{\lim}$ finies qui proviennent de $C$

[\uline{NB} l'inclusion $C\to\bar C$ commute aux $\varinjlim$ \deleted{\unsure{finies}} -- finies ou non --]

% loc: abelian variant
\deleted{\unsure{Ou bien peut supposer que}}\inserted{Supposons} $C$ est \deleted{\unsure{additive}} abélienne, pour $\varphi=$ ou encore $\psi\to I_2^{\operatorname{Hom}(\psi,\varphi)}$ est-il exact à g.
\[\underline{\operatorname{Hom}}(C,\mathrm{Ens})\unsure{\to|{-}}(\mathrm{Ens})\]

% loc: middle, injectivity
$\deleted{\operatorname{Hom}(\psi_1\times\psi_2,\varphi)\unsure{\to|\simeq}}$ Il faudrait que $\varphi$ transforme \unsure{monos|mono} \deleted{\unsure{(fibres)|en}} \inserted{en} épi i.e. $\varphi$ est ``injectif'' :

% loc: extension triangle
\diagram{\begin{tikzcd}[ampersand replacement=\&]\psi\arrow[r,hook]\arrow[dr]\&\psi'\arrow[d,dashed,"\exists"]\\\&\varphi\end{tikzcd}}

% loc: middle consequence
. Cela implique (faire $\psi=\varphi,\ \psi'=\varphi+\rho$) que pour $\forall\,\rho\ \unsure{\in|:}\, C\to\mathrm{Ens}$, on a $\underline{\operatorname{Hom}}(\rho,\varphi)\ne\varnothing$

% loc: left, second triangle
\diagram{\begin{tikzcd}[ampersand replacement=\&]\varphi\arrow[r]\arrow[dr,equal]\&\varphi\amalg\psi\arrow[d,dashed]\\\&\varphi\end{tikzcd}}

% loc: contradiction
\deleted{\unsure{autrement|C'est vrai si}} \uline{Mais c'est tjrs faux} \deleted{sauf si $C$ \unsure{est vide|vide} \illegible{1}} car si on prend pour $\rho$ le foncteur constant de valeur $x$ non vide, on a $\operatorname{Hom}(\rho,\varphi)=\varnothing$ car

% loc: Hom at empty object
\[\operatorname{Hom}_{\mathrm{Ens}}(\rho(\varnothing_C),\varphi(\varnothing_C))=\operatorname{Hom}_{\mathrm{Ens}}(x,\varnothing)=\varnothing\quad !\]

% loc: second question, data
\uline{Question}\qquad Soient $C$ petite catégorie avec $\varprojlim$ finies

$E$ catégorie avec $\varprojlim$ finies et petites $\varinjlim$

$\varphi:C\to E$ un foncteur exact à gauche

$\bar\varphi:\hat C\to E$ son prolongement canonique (commutant aux \inserted{petites} $\varinjlim$)

% loc: second question
Sous quelles conditions d'exactitude sur $E$ peut-on dire que $\bar\varphi$ est aussi exact à g ? [\uline{NB} si $E=(\mathrm{Ens})^\circ$, ce n'est \uline{jamais} le cas !]

% loc: last paragraph
On sait que c'est OK si $E$ est un topos. \uline{Réciproque} ? Supposons que $\varphi$ soit strictement générateur, \deleted{Alors le fait que} alors $E\to\hat C$ pl. \deleted{Il suffit que ce soit un topos.} fidèle et $\bar\varphi:\hat C\to E$ en est un adjoint à g. Donc si $\bar\varphi$ est exact à g, $E$ est un topos.

% loc: bottom right, upside down
% margin note written upside down
\diagram{\begin{tikzcd}[ampersand replacement=\&] \deleted{\unsure{C}/F}\arrow[d,dash]\&\&\\ \deleted{\unsure{C}}\&\deleted{x}\arrow[r]\arrow[d]\&\deleted{F}\\\&\deleted{y}\arrow[ur]\&\end{tikzcd}}

% archive sheet number: 15
